\subsection{Cross-Foundation Checks: Hybrid Pipeline and Full-Book Weighting}\label{sec:robustness-hybrid}

The hybrid pipeline of Section~\ref{sec:abm-danish} holds the paper's accounting fixed and swaps the micro-foundation for loan behavior. The ABM's macro-accounting layer runs unchanged (SOMA balances, phased-cap netting, income-scaled curtailment, term-aware amortization, the Danish dynamic-balance loop) while Path B's microsimulated CPR path replaces its CPR surface. No settlement kernel applies, since Path B settles prepayments in the month they occur. The U.S.-side recovery is 97.9\%, and the shared layer adds the curtailment flow Path B's standalone scorer omits: \$69.56 billion over the window. Under the production U.S.-intercept anchor, the Berger-recalibrated Danish re-run yields the U.S. 4.76\% / Danish 5.61\% mean-CPR pair and the $+\$61.2$ billion institutional gap Section~\ref{sec:abm-danish} adopts as its production estimate. The Danish-level bracketing anchor gives 3.39\% and $-\$99.9$ billion. One accounting layer with one empirically grounded micro-foundation carries both regimes, so accounting differences cannot confound the institutional comparison.

The same layer scores every estimator, which makes the recoveries commensurable with the benchmark. The actual SOMA roll-off behind the empirical shortfall contains borrowers' curtailment dollars, which no hazard simulation models. The shared layer adds them to each simulated roll-off. The netting is identical across estimators, \$69.56 billion for every run, and identical \emph{by construction}, not by coincidence. Hazard dynamics evolve on each estimator's own balances (Path A cohort balances, Path B loan balances, the Danish legs a deliberately compounded dynamic counterfactual, Appendix~\ref{sec:method-abm}). But every \emph{U.S.-leg} dollar conversion is common --- roll-off rescaled month-by-month to the actual holdings path (WSHOMCB) by the book-to-sample ratio. So the curtailment flow, computed on those holdings, is common, not per-run. It is a flat 9.1-percentage-point wedge between the two bases. The invariance covers U.S. legs only. The Danish legs curtail on their own counterfactual balances, whose slower drain retains more curtailment dollars. I quantify and print that differential: \$70.33 billion against the U.S. legs' common \$69.56 billion for the central hybrid (+\$0.77 billion). It is already embedded in every printed gap and immaterial to its sign or size.

Table~\ref{tab:bases} collects every hazard leg's recovery on all four accounting bases; the lock-in margin is basis-invariant at either floor and at every cell of Section~\ref{sec:robustness-floor}'s calibration box. The netting is bit-identical across legs and cancels exactly. The fully corrected figure comes from a joint run: the book-coupon-reweighted simulation scored through the shared layer in one pass. It recovers \$765.1 billion (unrounded: \$765.077 billion), 100.0\% of the \$764.7 billion benchmark at one decimal, exactly 100.04\% as the unrounded ratio $765.077/764.748$. That is a mixed-basis measurement. The reweight matches note rates to pass-through buckets. Converting the sample side to the pass-through basis first (note rate less the vintage guarantee fee and base servicing, $\Delta = 0.80$ points, legs at $0.60$--$0.90$) puts the like-for-like composed figure at 98.1\% (97.6--99.1 across the legs; run \texttt{coupon\_convention\_reweight}) --- so 100.0\% is the committed mixed-basis reading, not the corrected number. Path A's joint run reads \$905.1 billion, 118.4\%. The composed figure is conditional on the ex-ante covariate priors (92.2\%--99.2\% on the shared basis, approximately 94.3\%--101.3\% composed, Section~\ref{sec:pathb}). Its precision is therefore a property of that calibration, not a tuned result. Because the netting is common, the joint results equal the composed one-at-a-time deltas ($107.0 + 2.1 - 9.1$) exactly. That equality is a cross-check on the joint run, not its source. The abstract and Sections~\ref{sec:hazard-interp} and~\ref{sec:conclusion} state recoveries on this basis.

The curtailment layer's assumed rate is an external input, not a fitted quantity, so the layer invites one further stress. The monthly flow is an income-driven fraction, independent of that rate, times the rate on the actual holdings path (WSHOMCB). The netted total is therefore exactly proportional to the rate, and interacts with no estimator's dynamics. A $\pm 25\%$ stress (realized window average approximately 0.84\% annualized, Section~\ref{sec:method-benchmark}) scales the \$69.56 billion netting to \$52.17 and \$86.95 billion, and the 9.1-point wedge to 6.8 and 11.4 points. It shifts every U.S. leg's shared-basis recovery uniformly. The lock-in marginal and every cross-estimator gap stay unchanged to the decimal, since both legs of each comparison net the identical flow. Because the identity holds month by month, any time-varying shape subtracts out of the marginal term by term. Only the common level wedge and the Danish leg's sub-\$1 billion differential can move. Demonstration runs confirm it: seasonal, regime-split, and servicer-mix profiles all leave the marginal unchanged to below $10^{-13}$ points (Appendix~\ref{app:params}). The Danish legs, curtailing on their own counterfactual balances, scale only to first order. An exact accounting-layer rerun (CPR paths bit-identical across scales, curtailment entering only that layer) puts their differential at $+\$0.76$, $+\$0.77$, and $+\$0.66$ billion at $0.75\times$, $1.0\times$, and $1.25\times$. Those figures sit on the \emph{Danish-level bracketing} leg, whose anchor built the committed microsim cache the scaling run consumes, and they do not carry over. Re-scored on the U.S.-intercept cache (run \texttt{curtailment\_danish\_scaling}, production anchor), the differential reverses sign and crosses the materiality threshold: $-\$1.68$ billion at the unit point, $-\$2.39$ billion at $1.25\times$. Curtailment accrues on retained balance, so the sign follows which leg drains faster. The bracketing Danish leg is \emph{slower} than the U.S.\ leg (3.39\% against 4.76\% mean CPR) and accrues more; the production leg is \emph{faster} (5.61\%) and accrues less. I therefore withdraw the claim that the curtailment differential is immaterial to the Danish comparison. At $-\$1.68$ billion it is 2.7\% of the $+\$61.2$ billion in-sample gap and 6.0\% of the $+\$28.2$ billion off-window gap---small enough not to threaten either sign, but large enough to belong in the figure.

The second check reweights the sample to the book. Both headline recoveries are balance-weighted on the composition of the Freddie object each path uses: Path A the full 2017--2021 origination universe, Path B the 75,000-loan draw. That draw's unweighted mean note rate of 3.86\% sits above the book's 2.49\% face-weighted pass-through coupon, and it is where the Path B reweight starts. The unindented bucket shares Table~\ref{tab:composition} prints beside the book are the universe's, exposure-weighted at the window open --- different populations. Reweighting to the book's actual coupon mix shifts Path B from 107.0\% to 109.1\% and Path A from 121.5\% to 127.5\% on the committed mixed coupon basis. Those roughly two percentage points for Path B bound the sample-to-book extrapolation error along the coupon dimension. On a single pass-through convention the same reweight moves Path B only from 107.0\% to 107.2\%. The like-for-like composition effect is thus $+0.2$ points and the bound tightens by an order of magnitude (run \texttt{coupon\_convention\_reweight}; full sample-versus-book comparison in Appendix~\ref{app:composition}). The agency dimension (Ginnie Mae collateral) has no internal bound. It is externally bounded and now scored by a published-series overlay in Section~\ref{sec:pathb}, whose Ginnie-specific contribution ($+\$38.5$ billion) lands inside the static bound.
